\section{Introduction}

A game that averages 60 frames per second can still feel broken. Players
perceive a single 40\,ms frame in a stream of 16\,ms frames as a
\emph{hitch}, and reviewers and platform owners increasingly report the
1\% or 0.1\% low frame rate rather than the mean for exactly this reason.
On Linux the problem has grown with the platform: Proton~\cite{proton}
runs Windows games through Wine and a Direct3D-to-Vulkan translation
layer~\cite{dxvk,vkd3d-proton}, and handheld PCs such as the Steam Deck
run them on four cores and 16\,GB of memory shared between CPU and GPU. On
such machines a frame can miss its budget without the game doing anything
wrong: a lock holder is preempted, an allocation enters direct reclaim, an
asset read waits on slow storage, or a background process takes the CPU for
a few milliseconds.

These hitches are hard to diagnose because they are rare and unrepeatable.
A scheduler developer working on the Steam Deck's scheduler described
problems that appear as a spike every few minutes, with no way to trigger
tracing on the spike itself; he fell back to sampling every ten seconds and
named the missing piece as the correlation between the frame-level view and
the scheduler-level view~\cite{lwn-lpc2025-gaming,lpc2025-min}. In the same
session he was asked whether anyone had analyzed the chain of tasks from
the first wakeup to the last to explain a frame; nobody had.
\todo{Verify both statements against the LPC 2025 recording and quote it,
not LWN.}
Tracing everything for hours is not an answer: a full scheduler trace grows
faster than it can be stored, post-processing it takes hours~\cite{vapormark},
and the tracer perturbs the frame timing it is meant to explain.

The standard answer to ``capture only what matters'' is the \emph{flight
recorder}. The kernel keeps the last few seconds of events in an
overwriting ring buffer and freezes it when a trigger fires: an ftrace
snapshot trigger on a synthetic frame-latency event~\cite{ftracehist},
Perfetto's \texttt{STOP\_TRACING} trigger~\cite{perfetto-triggers}, LTTng
snapshots driven by latency-tracker~\cite{desfossez2016latency}, or
retroactive sampling in distributed systems~\cite{zhang2023hindsight}.
Flight recorders solve capture, and we do not claim otherwise. They leave
three problems. First, the trigger sees only the symptom: a frame that is
late because the game spent 25\,ms computing pathfinding fires the same
trigger as a frame that is late because the render thread waited on a
preempted worker. Second, each trigger dumps the whole window, megabytes of
raw events, so output grows with the number of hitches times the window
size, and a game that hitches every few seconds over a multi-hour session
brings back the storage problem. Third, the dump is raw: an engineer still
has to find, among thousands of scheduler events, which thread the frame
was waiting on and what that thread was waiting on in turn.

Per-event thresholds (\texttt{perf --off-cpu-thresh}~\cite{perf-offcpu-thresh},
bcc's \texttt{*slower} tools~\cite{gregg2019bpf}, ftrace filters) and kernel
delay accounting~\cite{delayacct} address output volume, but they decide at
the wrong time. A per-event filter judges each stall when it ends, with only
that stall in hand. It cannot keep a 400\,$\mu$s futex wait because the frame
it belongs to will turn out to be 30\,ms long, so a hitch made of forty
sub-threshold stalls is invisible to it. Delay accounting accumulates without
a threshold, but has no notion of a frame, no futex category and no
waker.

Our observation is that what should cross into user space is neither the
raw events nor the individual stalls, but an \emph{attribution record} for
the frame: how the presenting thread's frame time divides among kernel
causes, and which chain of other threads and interrupts its longest wait
depended on. That record can only be computed from state accumulated across
the whole frame and must be committed or discarded when the frame ends.
Both require stateful, programmable logic that runs at the event, which is
what eBPF provides and what ftrace's declarative filter, trigger and
histogram language deliberately does not (\autoref{sec:ftrace}).

This paper presents \sys. \sys has three parts:
\begin{itemize}
\item \textbf{A frame source that works for unmodified games.} An implicit
  Vulkan layer brackets \texttt{vkQueuePresentKHR} with two exported marker
  functions that \sys probes. A uprobe on the Vulkan loader's exported
  symbol would see nothing: applications that fetch entry points through
  \texttt{vkGetDeviceProcAddr} (volk, DXVK, vkd3d-proton, Wine's Vulkan
  bridge, and \texttt{vkcube} itself) bypass the loader
  entirely~\cite{vulkan-loader}.
\item \textbf{A frame-scoped partition of the presenting thread's time.}
  Every interval of the root thread within the frame is clipped to the
  frame window and assigned to exactly one of fifteen cause buckets, so the
  buckets sum to the frame time. This is maintained in the kernel, in
  task-local storage, as the scheduler runs.
\item \textbf{Online wake-chain resolution.} Every task keeps a short
  history of its own off-CPU time and carries a link to its waker: how long
  the waker was itself off-CPU while it was being waited for, who last held
  the waker's CPU, and, recursively, the waker's own links. When the root's
  blocked interval ends, the time its waker was itself off-CPU is
  \emph{inherited} (a kernel stall one hop away) and the remainder is
  \emph{waiting on a running task} (the game waiting for its own work, not
  a kernel stall). The root's waits are grouped by cause and by who ended
  them, each group with its chain. No graph is walked at report time.
\end{itemize}
A record crosses into user space only when the frame exceeds its budget
and, optionally, only when a configurable share of the overrun is
kernel-observable.

We implement \sys in 1,828 lines of BPF C and 1,992 lines of user-space C,
reusing a symbolizer and interrupt-context detection from our earlier
wake-chain profiler, plus a 689-line Vulkan layer. We also build \bench, a
synthetic frame loop with \hbnClasses\ injectable stall classes and a
per-frame ground-truth log. On \bench, \sys recorded all \hbtotOver\
over-budget frames of \hbtotFrames\ and no frame within budget; its largest
bucket named the injected mechanism in \hbtotMatch\ of \hbtotInj\ injected
frames, with the misses explained in \autoref{sec:eval-accuracy}; chain
resolution cleanly separated a thread waiting on a stalled worker from one
waiting on a busy worker; and a gate that uses the resolution rejected every
negative control while keeping every positive class. We then reproduce ten
mechanisms behind documented Linux gaming stutters around an unmodified
game-shaped program: the records explain nearly every affected frame, and
three of the cases found faults in an earlier version that \bench had not.
On the same hitches made rare, \sys is the only one of eight tools that
attributes nearly every hitch and names its culprit, and the raw recorders
that also hold the frames, Perfetto's flight recorder among them, need 155
to 900 times its output per hitch. Its
overhead per context switch is in the range of existing scheduler tracers,
not below it; what it saves is output.
\tbd{[One sentence on Proton titles, one on prevalence.]}

This paper makes the following contributions.
\begin{itemize}
\item We identify the \emph{symptom--cause gap} in hitch capture and show
  which parts of frame-scoped attribution an in-kernel declarative tracer
  can and cannot express, locating exactly where programmable in-kernel
  logic is needed (\autoref{sec:problem}).
\item We present a design that partitions a frame's critical thread by kernel
  cause and resolves its waits through wake chains online, with a commit or
  discard decision at frame end, so that output is proportional to hitches
  and normal frames emit nothing (\autoref{sec:design}).
\item We show how to express this under the eBPF verifier on a stock
  kernel, including the restructurings the verifier forced
  (\autoref{sec:verifier}).
\item We build \bench, a ground-truth hitch benchmark, and ten case
  studies that reproduce documented stutter mechanisms, and evaluate
  attribution, chain resolution, cause gating, how eight tools capture and
  explain rare hitches, and overhead, reporting what the records cannot yet
  say; \tbd{prevalence on real titles} (\autoref{sec:eval}).
\end{itemize}
